\documentclass[twocolumn,11pt,a4paper]{article}

\usepackage[utf8]{inputenc}
\usepackage[T1]{fontenc}
\usepackage{lmodern}
\usepackage[margin=1in]{geometry}
\usepackage{amsmath, amssymb, amsthm, mathtools}
\usepackage{thmtools}
\usepackage{enumitem}
\usepackage{microtype}
\usepackage{booktabs}
\usepackage{graphicx}
\usepackage{caption}
\usepackage{subcaption}
\usepackage{hyperref}
\usepackage{cleveref}
\usepackage{xcolor}
\usepackage{tikz}
\usetikzlibrary{arrows.meta,positioning,calc}

\hypersetup{
  colorlinks=true,
  linkcolor=blue!50!black,
  citecolor=green!50!black,
  urlcolor=blue!60!black
}

%---- theorem environments ----
\declaretheoremstyle[
  spaceabove=6pt, spacebelow=6pt,
  headfont=\bfseries, notefont=\mdseries, notebraces={(}{)},
  bodyfont=\itshape, postheadspace=1em
]{thmstyle}

\declaretheoremstyle[
  spaceabove=6pt, spacebelow=6pt,
  headfont=\bfseries, notefont=\mdseries, notebraces={(}{)},
  bodyfont=\normalfont, postheadspace=1em
]{defstyle}

\declaretheorem[style=thmstyle, name=Theorem, numberwithin=section]{theorem}
\declaretheorem[style=thmstyle, name=Proposition, sibling=theorem]{proposition}
\declaretheorem[style=thmstyle, name=Lemma, sibling=theorem]{lemma}
\declaretheorem[style=thmstyle, name=Corollary, sibling=theorem]{corollary}
\declaretheorem[style=defstyle, name=Definition, sibling=theorem]{definition}
\declaretheorem[style=defstyle, name=Axiom, sibling=theorem]{axiom}
\declaretheorem[style=defstyle, name=Remark, sibling=theorem]{remark}
\declaretheorem[style=defstyle, name=Example, sibling=theorem]{example}
\declaretheorem[style=defstyle, name=Principle, sibling=theorem]{principle}

%---- notation macros ----
\newcommand{\R}{\mathbb{R}}
\newcommand{\Q}{\mathbb{Q}}
\newcommand{\Z}{\mathbb{Z}}
\newcommand{\N}{\mathbb{N}}
\newcommand{\C}{\mathbb{C}}
\newcommand{\F}{\mathcal{F}}
\newcommand{\K}{\mathcal{K}}
\newcommand{\D}{\mathcal{D}}
\newcommand{\M}{\mathcal{M}}
\newcommand{\X}{\mathcal{X}}
\newcommand{\Y}{\mathcal{Y}}
\newcommand{\W}{\mathcal{W}}
\newcommand{\Cell}{\mathcal{C}}
\newcommand{\Mode}{\mathfrak{M}}
\newcommand{\Method}{\mathfrak{P}}
\newcommand{\receiver}{\mathsf{R}}
\newcommand{\Sfloor}{S_{\flat}}
\newcommand{\eps}{\varepsilon}
\newcommand{\norm}[1]{\lVert#1\rVert}
\newcommand{\abs}[1]{\lvert#1\rvert}
\newcommand{\inner}[2]{\langle #1,#2\rangle}
\newcommand{\Span}{\mathrm{span}}
\newcommand{\diam}{\mathrm{diam}}
\newcommand{\dom}{\mathrm{dom}}
\newcommand{\codom}{\mathrm{codom}}
\newcommand{\id}{\mathrm{id}}
\newcommand{\dd}{\,\mathrm{d}}

\title{\textbf{Constituents of Charge Dynamics in Bounded Microfluidic Networks:\\
Sufficiency, Closure, and Path-Independent Navigation in Multi-Receiver Systems}}

\author{Kundai Farai Sachikonye\\[2pt]
\small Technical University of Munich \\
\small Chair of Proteomics and Bioanalytics, Experimental Bioinformatics\\
\small Maximus-von-Imhof-Forum 3, 85354 Freising, Germany\\
\small \texttt{kundai.sachikonye@wzw.tum.de}}
\date{\today}

\begin{document}
\maketitle

\begin{abstract}
We develop a unified framework for understanding charge circulation in bounded microfluidic networks with multiple receivers and feedback architectures. The principal result is a theory of how navigation to action-cells (behavioral termini) is achieved through path-independent convergence, despite unbounded variability in internal trajectory structure. We establish five theorem clusters: (i)~\emph{Sufficiency via Unconstrained Subtasks}: global viability of charge circulation is maintained despite arbitrary subtask variation, provided the S-functional (residual distance to goal) remains within bounds; (ii)~\emph{Topological Closure Requirement}: systems that generate outbound charge must maintain inbound return paths for circuit completion, independent of signal fidelity; (iii)~\emph{Poincar\'{e} Deviation in Feedback}: return signals are asymptotically close to sent commands but never exact, inducing strictly positive deviation that generates successive distinct states; (iv)~\emph{Fabrication with Correction}: systems generate internal charge representations (fabrications) that are corrected by external feedback when inputs are available, and unconstrained when inputs vanish; (v)~\emph{Production--Completion Incompatibility}: knowledge production (dispersion $\sigma>0$ in methodology) and action completion (commitment at $\sigma=0$) are mutually exclusive at each iteration. Three equivalent representations (oscillatory, categorical, partition) make all claims representation-invariant. We prove that path-independent navigation is possible when multiple approach trajectories satisfy the sufficiency condition, independent of their intermediate structure. We show that systems lacking feedback paths cannot maintain stable navigation despite identical outbound commands. We establish that fabrication amplitude is unconstrained by internal consistency alone and requires external correction. Forty numerical experiments validate all principal theorems to machine precision. The framework applies to any bounded system with mobile charges in confined environments, providing mathematical foundations for understanding navigation, action-commitment, and feedback requirements in networks of receivers.
\end{abstract}

\section{Introduction}
\label{sec:introduction}

%==========================================================
\subsection{Motivation: Path Independence in Bounded Systems}

Consider a microfluidic network where multiple independent pathways must navigate to an identical action-cell (a designated region of configuration space representing successful task completion). Three pathways arrive at this cell:

\begin{itemize}[leftmargin=*]
\item Pathway A enters with high initial charge concentration, rapid transients, noisy intermediate states
\item Pathway B enters with low initial charge, monotonic progression, clean transients
\item Pathway C couples to external pathways, its trajectory shaped by neighbor-state feedback
\end{itemize}

Despite these radically different approach trajectories, all three reach the identical action-cell terminus. By what principle? Standard control theory would demand trajectory-by-trajectory specification—each path individually tuned to reach the goal. Yet microfluidic and biological systems achieve identical outcomes through diverse routes.

This paper establishes that the answer lies in \emph{sufficiency}: as long as the global S-functional (residual distance to the action-cell) remains bounded within a critical threshold, the specific trajectory is irrelevant. The system can take any path, through any intermediate state, with any intermediate charge distribution—provided the final destination is reached.

This principle—path independence through sufficiency—is the mathematical core of navigation in bounded systems.

%==========================================================
\subsection{The Closed-Loop Requirement}

A second observation: systems that generate outbound charge command must maintain inbound return paths. This is not a control requirement; it is a topological requirement.

Consider a motor pathway that sends outbound charge (a command signal) into a mechanical system. If this outbound charge has no return path—no feedback loop, no inbound signal to close the circuit—then something remarkable happens: the network cannot generate stable navigation, even though the outbound command is identical to cases where return paths exist.

This is the \emph{circuit closure theorem}: movement and stable action require topologically closed charge circulation. The circuit must be geometrically closed (charge out, charge back). This closure is necessary not for accurate feedback (the return signal can be noisy, delayed, or partially fabricated) but for maintaining the circuit topology itself.

%==========================================================
\subsection{Fabrication and Correction}

A third principle emerges from systems operating in different input regimes:

\emph{With external input available}, the system generates internal charge representations (a model of the expected state) and corrects these representations against actual input. When external input contradicts the internal model, the model is updated.

\emph{Without external input}, the system continues to generate internal charge representations unconstrained. There is no external signal to correct against. The generation proceeds freely, producing ever-more-varied internal states.

This distinction explains a paradoxical finding: systems can maintain fully functional operation while generating increasingly inaccurate internal representations, \emph{provided the circuit topology remains closed}. The fidelity of the internal model is decoupled from the system's ability to maintain circuit function.

%==========================================================
\subsection{Production and Completion as Incompatible Phases}

A fourth principle establishes an incompatibility. Consider two operational regimes:

\begin{enumerate}[label=(\Roman*),leftmargin=*]
\item \textbf{Knowledge Production}: the system generates multiple possible outcomes for a given input (dispersion $\sigma>0$ in the methodology), producing variation that explores possibility space.

\item \textbf{Action Completion}: the system commits to a single outcome and executes it deterministically ($\sigma=0$ at the decision step).
\end{enumerate}

These regimes are mutually exclusive at each iteration. A system cannot simultaneously explore multiple possibilities and commit to a single action. It must alternate: produce (explore), complete (commit), produce, complete.

This incompatibility has profound consequences. It means that any system performing navigation must have two distinct operational modes, activated alternately. The implications for receiver architecture and feedback design are substantial.

%==========================================================
\subsection{Organization and Scope}

This paper develops a mathematical framework for bounded microfluidic networks integrating these four principles: sufficiency, closure, fabrication-with-correction, and production--completion incompatibility.

We prove:
\begin{enumerate}[label=(\arabic*),leftmargin=*]
\item Path-independent convergence to action-cells (Theorem 1)
\item Topological closure as necessity (Theorem 2)
\item Poincar\'{e} deviation in feedback systems (Theorem 3)
\item Fabrication bounds and correction requirements (Theorem 4)
\item Incompatibility of production and completion (Theorem 5)
\item Equivalence across three representations (Theorem 6)
\end{enumerate}

We establish that these principles are interconnected through the S-functional framework, where $S(\receiver, x; \Cell)$ measures the residual distance from state $x$ to an action-cell $\Cell$ relative to a receiver $\receiver$. The sufficiency principle states that if $S$ is below a threshold (the action-cell's tolerance), the approach trajectory is operationally irrelevant.

The framework applies to any bounded system with mobile charges in confined geometries. Specific applications include genetic networks, metabolic pathways, neural microarchitectures, and distributed sensing systems. However, we present the theory in abstract form, allowing the reader to recognize applications specific to their domain.

\section{Foundations and Notation}
\label{sec:foundations}

[Content to follow in subsequent sections...]

\bibliographystyle{plain}
\bibliography{references}

\end{document}
